% Part: incompleteness
% Chapter: theories-computability
% Section: consis-with-q

\documentclass[../../../include/open-logic-section]{subfiles}

\begin{document}

\olfileid{inc}{tcp}{con}

\olsection{$\Th{Q}$-এর সঙ্গে সঙ্গতিপূর্ণ তত্ত্বগুলি অনির্ণেয়}

পরের উপপাদ্যটি বলে যে শুধু $\Th{Q}$-ই অনির্ণেয় নয়; বস্তুত,
$\Th{Q}$-এর সঙ্গে বিরোধ নেই এমন যেকোনো তত্ত্বই অনির্ণেয়।

\begin{thm}
ধরি $\Th{T}$ পাটিগণিতের ভাষার এমন যেকোনো তত্ত্ব যা $\Th{Q}$-এর সঙ্গে
সঙ্গতিপূর্ণ (অর্থাৎ, $\Th{T} \cup \Th{Q}$ সঙ্গতিপূর্ণ)। তাহলে
$\Th{T}$ অনির্ণেয়।
\end{thm}

\begin{proof}
মনে করো, $\Th{Q}$-এর স্বতঃসিদ্ধের সসীম সেটটি হলো $!Q_1$, \dots,
$!Q_8$। এমনকি এগুলিকে একটি স্বতঃসিদ্ধ
$!E = !Q_1 \land \dots \land !Q_8$ দিয়ে প্রতিস্থাপন করা যায়।

ধরি $\Th{T}$ একটি নির্ণেয় তত্ত্ব এবং $\Th{Q}$-এর সঙ্গে সঙ্গতিপূর্ণ।
নিই
\[
C = \Setabs{!A}{\Th{T} \Proves !E \lif !A}.
\]
দেখাই যে $C$ তাহলে $\Th{Q}$ ও $\Th{\bar Q}$-এর একটি গণনসাধ্য
পৃথকীকরণ হবে, যা বিরোধ। প্রথমত, $!A$ যদি $\Th{Q}$-তে থাকে, তাহলে
$!A$ $\Th{Q}$-এর স্বতঃসিদ্ধগুলি থেকে প্রমাণযোগ্য; নিঃসরণ উপপাদ্য
অনুসারে প্রথম-ক্রমের যুক্তিবিদ্যায় $!E \lif !A$-এর কোনো
!!a{derivation} আছে। তাই $!A$ $C$-তে আছে।

অন্যদিকে, $!A$ যদি $\Th{\bar Q}$-তে থাকে, তাহলে প্রথম-ক্রমের
যুক্তিবিদ্যায় $!E \lif \lnot !A$-এর একটি প্রমাণ আছে। $\Th{T}$ যদি
$!E \lif !A$-ও প্রমাণ করে, তবে $\Th{T}$ $\lnot !E$ প্রমাণ করে; সে
ক্ষেত্রে $\Th{T} \cup \Th{Q}$ অসঙ্গতিপূর্ণ। কিন্তু আমরা
$\Th{T} \cup \Th{Q}$-কে সঙ্গতিপূর্ণ ধরেছি; তাই $\Th{T}$
$!E \lif !A$ প্রমাণ করে না, এবং $!A$ $C$-তে নেই।

আমরা দেখিয়েছি, $!A$ $\Th{Q}$-তে থাকলে সেটি $C$-তে থাকে, আর $!A$
$\Th{\bar Q}$-তে থাকলে সেটি $\Complement{C}$-তে থাকে। অতএব $C$ একটি
গণনসাধ্য পৃথকীকরণ; এটাই অভীষ্ট বিরোধ।
\end{proof}

এই উপপাদ্যটি খুব শক্তিশালী। যেমন, এটি থেকে পাওয়া যায়:

\begin{cor}
  পাটিগণিতের ভাষার প্রথম-ক্রমের যুক্তিবিদ্যা (অর্থাৎ সেটটি
  $\Setabs{!A}{\text{$!A$ প্রথম-ক্রমের যুক্তিবিদ্যায় প্রমাণযোগ্য}}$)
  অনির্ণেয়।
\end{cor}

\begin{proof}
প্রথম-ক্রমের যুক্তিবিদ্যা হলো সেই $\emptyset$-এর অনুগমনগুলির সেট যা
$\Th{Q}$-এর সঙ্গে সঙ্গতিপূর্ণ।
\end{proof}

\end{document}
